% PASS20260921 twin E8 order-six Kac classification.
\subsection*{The two physicalized order-six actions are distinct affine-$E_8$ Kac classes}

The structural FI $\times$ matter-parity quotient and the flagship physical
holonomy both define certified inner order-six actions on the $E_8$ adjoint,
but their eigenspace fingerprints differ:
\[
(54,48,33,32,33,48)
\qquad\hbox{versus}\qquad
(44,40,38,48,38,40).
\]
Using the affine highest-root marks in the repository's frozen simple-root
ordering,
\[
(a_0,\ldots,a_8)=(1,4,6,5,4,3,2,2,3),
\]
an exhaustive Cartan--Kac classification of all normalized exact order-six
inner diagrams finds exactly twenty candidates and selects uniquely
\[
\boxed{s_{\rm FI\times P_M}=(2,0,0,0,1,0,0,0,0)}
\]
and
\[
\boxed{s_{\rm hol}=(0,0,0,0,1,0,0,1,0)}.
\]
Their power ladders are therefore
\[
\begin{array}{c|ccc}
 &g&g^2&g^3\\ \hline
{\rm FI}\times P_M&
D_5+A_2+\mathfrak u_1&E_6+A_2&D_8\\
{\rm holonomy}&
D_4+A_3+\mathfrak u_1&D_7+\mathfrak u_1&D_8 .
\end{array}
\]
The fixed dimensions are respectively
\[
(54,86,120)
\qquad\hbox{and}\qquad
(44,92,120).
\]
Thus the previously certified common $D_8$ cube is the exact intersection of
the two power ladders, while the square already separates the conjugacy
classes.  Equivalently, both Kac diagrams use the same mark-four finite node;
the remaining Kac weight two is carried by the affine node in the structural
class and by a finite mark-two node in the holonomy class.

This is the standard Cartan--Kac torsion classification (see M.~Reeder,
\emph{Enseign. Math.} 56 (2010), 3--47).  The executable witness is
\texttt{analysis/w33\_e8\_order6\_kac\_classification.py}.

\paragraph{Scope firewall.}
This is an exact conjugacy-class theorem inside complex $E_8$.  It neither
proves a matter-parity-preserving heterotic D/F-flat vacuum nor identifies the
two distinct order-six classes with one another or with additional spacetime
orbifold data.
